\section{Related Works}

This section reviews major developments in text-to-video generation, categorized by proprietary systems, open-source efforts, and recent trends in autoregressive and causal modeling. We highlight unresolved challenges in scalability, causality, and streaming compatibility—challenges that \magi is designed to address.

\paragraph{Proprietary Systems.}
Recent proprietary models have significantly advanced generation length, resolution, and semantic fidelity. OpenAI’s Sora~\citep{openaisora2024} introduced long-form, high-resolution generation with strong prompt consistency. Kuaishou’s Kling~\citep{kuaishou2024kling} and Runway’s Gen‑3~\citep{runway2024gen3} emphasized temporal fidelity and fine-grained stylistic control, respectively. Luma AI’s DreamMachine~\citep{luma2024dm} improved motion continuity and stylistic adherence. Pika Labs’ Pika 1.5~\citep{pika2024pika} enabled interactive control over visual attributes, while Meta’s MovieGen~\citep{polyak2024moviegencastmedia} offered transparency into foundational model training. Most recently, Google's Veo 2~\citep{veo22025} advanced physical realism and human motion modeling. Despite these innovations, most systems are closed-source and opaque in architecture, limiting reproducibility and extensibility.

\paragraph{Open-Source Ecosystem.}
The open-source community pioneered latent diffusion through Stable Diffusion~\citep{esser2020taming}, which integrated a variational autoencoder~\citep{kingma2013auto} for latent representation, a CLIP-based text encoder~\citep{radford2021learning}, and a U-Net denoiser~\citep{ronneberger2015u}. Temporal extensions such as VDM~\citep{ho2022video}, AnimateDiff~\citep{guo2023animatediff}, and SVD~\citep{blattmann2023svd} adapted the architecture for frame coherence. Transformer-based backbones like DiT~\citep{dit}, PixArt‑$\alpha$~\citep{chen2023pixartalpha}, and Latte~\citep{ma2024latte} demonstrated scalability and inspired early video adaptations. Recent open implementations—including Open‑Sora~\citep{opensora}, Open‑Sora‑Plan~\citep{lin2024open}, CogVideoX~\citep{cogvideox}, Mochi 1~\citep{genmo2024mochi}, HunyuanVideo~\citep{kong2024hunyuanvideo}, StepVideo~\citep{ma2025step}, LTX‑Video~\citep{HaCohen2024LTXVideo}, and Wan~\citep{wang2025wan}—introduced modular advances in chunking, compression, and streaming. However, these systems largely retain bidirectional denoising and globally conditioned inference, limiting applicability to real-time or causal settings.

\paragraph{Autoregressive and Causal Modeling.}
An emerging trend is the integration of autoregressive modeling and causal constraints. Diffusion Forcing~\citep{chen2024diffusion} introduces independent per‑token noise schedules that allow a causal model to denoise future tokens while keeping past tokens minimally perturbed, effectively unifying next‑token prediction with full‑sequence diffusion. FVDM~\citep{liu2024redefining} employed timestep vectorization for precise noise control. CausVid~\citep{yin2024slow} combined causal inference with distillation for streaming scenarios. While promising, these models remain limited in scale, often lack chunk-wise abstraction, and do not unify video continuation with I2V/T2V generation.

\paragraph{\magi: Scalable Autoregressive Diffusion.}
To our knowledge, \magi is the first large-scale, chunk-wise autoregressive diffusion model trained from scratch that unifies high-fidelity text-to-video, image-to-video, and video continuation tasks under strict causal constraints. It supports real-time streaming and long-horizon synthesis via efficient chunk-wise denoising, shortcut distillation, and KV-cached inference. By explicitly addressing scalability, causality, and streaming compatibility, \magi establishes a new foundation for unified and controllable video generation.
